\documentclass[12pt]{article}
\usepackage{amsmath,amssymb}

\begin{document}
\section*{{2024 AIME II Problems/Problem 9}}
Source: AoPS Wiki

\subsection*{Solution}
The problem says "some", so not all cells must be occupied.
We start by doing casework on the column on the left. There can be 5,4,3,2, or 1 black chip. The same goes for white chips, so we will multiply by 2 at the end. There is $1$ way to select $5$ cells with black chips. Because of the 2nd condition, there can be no white, and the grid must be all black- $1$ way . There are $5$ ways to select 4 cells with black chips. We now consider the row that does not contain a black chip. The first cell must be blank, and the remaining $4$ cells have $2^4-1$ different ways, since each cell can be white or blank( $-1$ comes from all blank). This gives us $75$ ways. Notice that for 3,2 or 1 black chips on the left there is a pattern. Once the first blank row is chosen, the rest of the blank rows must be ordered similarly. For example, with 2 black chips on the left, there will be 3 blank rows. There are 15 ways for the first row to be chosen, and the following 2 rows must have the same order. Thus, The number of ways for 3,2,and 1 black chips is $10*15$ , $10*15$ , $5*15$ . Adding these up, we have $1+75+150+150+75 = 451$ . Multiplying this by 2, we get $\boxed{902}$ .
~westwoodmonster

\end{document}
